\begin{lemma}[Claim 9: first tail use at the seed]
Fix the Section 7 minimum context and a
\noderef{HierarchyReplacementPresentation} \(P\).  Let \(\alpha\) be
missing on \(P\)'s last starred stage \(A\) and at \(y_\ell\).  Assume
the ambient Situation 7.1 condition
\(\noderef{HierarchyPathNumber}
(P.\mathit{counterexample}.\mathit{hierarchy})=1\).  Assume
the normalized alternative that \(\alpha\) is missing on \(T_{n,q}\)
when \(q>0\), and on \(T_n\) when \(q=0\).  Suppose no
\(y_i\), \(1\leq i\leq\ell-1\), has a color common-missing with
\(y_\ell\) under \(P\)'s coloring, and assume every proper
exactly-one-uncolored replacement coloring has a missing color at every
vertex.  If the exact first edge after \(A\), at position \(|A|\), is
colored \(\alpha\), then
\noderef{HierarchyEligibleCommonMissingIndices} is nonempty.
The proof constructs the paper's closed color \(\gamma\), performs the
specified Kempe change while preserving the exact first edge, and reduces
to either the unused-tail case or the first-use-\(j\geq2\) case.
\end{lemma}
\begin{proof}
Write
\[
 \mu_n=P.\mathit{coloring},\quad
 H=P.\mathit{counterexample}.\mathit{hierarchy},\quad
 A=\noderef{HierarchyStarredStage}(H,q),\quad
 \ell=\noderef{HierarchyFullTailLength}(C),
\]
where \(q\) is the last stored index.  The
\noderef{SameGoodHierarchy} field of \(P\), unfolded through
\noderef{SameHierarchyStages}, identifies the original and current
divider counts, base core, stages, last starred stage, and hence all
full-tail prefixes.  Let
\[
 X=T(y_1)=\noderef{HierarchyFullTailPrefix}(C,1),\qquad
 B=T(y_{\ell-1})=\noderef{HierarchyFullTailPrefix}(C,\ell-1),
\]
and let \(y\) be the final displayed vertex of \(T\).
\noderef{HierarchyCounterexampleTailAtLeastTwo} and
\noderef{HierarchyCounterexamplePenultimateElementary}, applied to
\(P.\mathit{minimumContext}\), give
\[
 2\le\ell,\qquad
 \noderef{ColoredVertexSetElementary}
       (G,\mu_n,\noderef{TreeSequenceVertexSet}(B)).          \tag{1}
\]
The nested and terminal-stage fields of
\noderef{TashkinovHierarchy}, together with the exact prefix equations,
give
\[
 A\prec X\preccurlyeq B\prec T,\quad |X|=|A|+1,\quad
 e_1=T.\mathit{edge}(|A|).                                  \tag{2}
\]

We first choose the auxiliary color with the exclusion required by the
paper:
\[
 \gamma\in\overline\mu_n(A),\qquad
 \gamma\notin\bigcup_{\eta\in D_{n,q}}
       P.\mathit{counterexample}.\mathit{good}.\Gamma^q_\eta,
 \qquad \gamma\ne\alpha,                                    \tag{3}
\]
and \(\{\gamma\}\) is closed on \(A\).  Here are the three cases proving
existence.

Suppose \(q>0\).  Then \(A=T_{n,q}\).  The terminal-ladder missing set is
contained in the missing set of \(A\), since \(T_n\preccurlyeq A\), so
the context bound gives at least \(2n+11\) colors missing on \(A\).
\noderef{HierarchyStoredColorsCardinality} gives
\(|D_{n,q}|\le n\), and every Gamma pair has size two by
\noderef{GoodHierarchy}.twoColors.  At the last refresh,
\noderef{HierarchyStoredColorsAntitone} puts every currently stored color
in the preceding stored set, and
\noderef{GoodHierarchy}.uniqueRefresh says that all pairs except one are
copied from the preceding index; the exceptional new pair has size two.
Consequently the union of the preceding Gamma family and the full last
Gamma family has size at most \(2n+2\).  After also deleting the singleton
\(\{\alpha\}\), at least eight of the \(2n+11\) missing colors remain.
Choose \(\gamma\) among them.  The last
\noderef{GoodHierarchy}.minusClosed clause says that every color missing
on \(T_{n,q}=A\) outside the preceding Gamma union has empty boundary
there.  Thus (3) holds and \(\{\gamma\}\) is
\noderef{TreeColorSetClosed} on \(A\).

Suppose \(q=0\) and the final kind is SE or RE.  The nonparallel clause of
\noderef{HierarchyBase} gives \(A=T_n\).  Again the context gives at
least \(2n+11\) missing colors, while the zeroth Gamma union has size at
most \(2n\); deleting it and \(\{\alpha\}\) leaves a choice for
\(\gamma\).  We prove that every color missing on \(T_n\) is closed under
\(\mu_n\).  Inspect the actual final
\noderef{TashkinovSeriesStep}.  In RE, its current-coloring equality and
preceding \noderef{ColoredVertexSetClosed} field give this immediately.
In SE, its \noderef{MaximumDefectiveChoice}.selectedStable field and the
symmetry theorem \noderef{sc2} make the preceding coloring stable relative
to \(\mu_n\).  If a \(\mu_n\)-colored boundary edge had a color missing
on \(T_n\), the reverse incident-edge stability clause would give the
same edge the same color under the preceding coloring, while vertexwise
stability would preserve the missing witness.  This contradicts the
preceding closedness.  Hence the selected \(\gamma\) is singleton-closed
on \(A=T_n\).

Finally suppose \(q=0\) and the final kind is PE.  The
\noderef{GoodHierarchy}.parallelBase cardinality clause gives at least
four colors in
\[
 \bigl(\overline\mu_n(T_n)\cap
       \overline\mu_n(P.\mathit{replacementBase}.\mathit{closure})\bigr)
       \setminus\Gamma^0.
\]
After excluding \(\alpha\), choose \(\gamma\) from the at least three
remaining colors.  To establish closedness, use actual predecessor
elementarity.  If the rung is one, it is the first conclusion of
\noderef{TechnicalTheoremBase}.  If the rung is greater than one,
\noderef{CertifiedETTRestriction} at rung \(n-1\) supplies an actual
maximum predecessor ETT on \(T_n\), and the lower-rung elementarity field
of the minimum context applies.  Therefore
\noderef{ParallelCommonMissingClosed} says that every color in the
displayed intersection, including \(\gamma\), is closed on the exact PE
ordered-union core \(A\).  This completes the proof of (3) in all cases.
In particular, \(\gamma\ne\alpha\) follows from the deliberate singleton
exclusion, not from incidence of \(e_1\).

We apply \noderef{ChangeLemma} with the closed color \(\gamma\) as its
first color and \(\alpha\) as its second color.  Both colors are missing
on \(A\preccurlyeq B\), so the conditional order-and-unused-edge premise
is vacuous.  Its elementarity premise is exactly (1), whose
pairwise-disjoint missing-set clause is the definition of
\noderef{ColoredVertexSetElementary}.  The third, closed-color conclusion
gives
\[
 K_A:=\noderef{TwoColorComponent}(G,\mu_n,\gamma,\alpha,
            \operatorname{inl}v_\gamma)
     =\noderef{TwoColorComponent}(G,\mu_n,\gamma,\alpha,
            \operatorname{inl}v_\alpha),                    \tag{4}
\]
and says that \(\gamma,\alpha\) are interchangeable on \(A\), including
possibly singleton path components.

Let \(Q\) be the full \((\gamma,\alpha)\)-component through
\(\operatorname{inl}y\).  It is a
\noderef{TwoColorPathComponent}, because \(\alpha\) is missing at \(y\),
properness gives incidence degree at most one there by
\noderef{MissingColorIncidenceEndpoint}, and
\noderef{TwoColorComponentSpanningPath} spans its component by a simple
path.  We claim \(Q\) does not meet \(V(A)\).  Otherwise interchangeability
from (4) identifies \(Q\) with the path component \(K_A\) meeting \(A\).
The latter is indeed a path component: its root \(v_\gamma\) is missing
\(\gamma\), so \noderef{MissingColorIncidenceEndpoint} bounds the root
degree by one and \noderef{TwoColorComponentSpanningPath} supplies a
simple spanning path.
If \(v_\gamma=v_\alpha\), both colors are missing at that one vertex, so
its incidence node is isolated and its component cannot also contain the
distinct vertex \(y\).  If they are distinct, the two missing vertices
have disjoint missing sets by the pairwise clause of
\noderef{ColoredVertexSetElementary} in (1), while the exact strict prefix
\(A\prec T\) and vertex injectivity make each of them distinct from
\(y\).  Thus they and \(y\) are three distinct incidence nodes of degree
at most one in the same component.  This contradicts
\noderef{TwoColorIncidenceDegreeAtMostTwo} and the two-ended spanning
path, exactly as in the unused-tail case.  Hence \(Q\cap V(A)=\varnothing\).
In particular \(Q\ne K_A\), and
\noderef{TwoColorComponentsDisjoint} upgrades this inequality to
disjointness of their full vertex-and-edge-node supports.
By adjacency in \noderef{TwoColorIncidenceGraph}, no edge-node of \(Q\)
can be incident with a vertex of \(A\).

The exact edge \(e_1\) is incident with its predecessor \(A\) by the
attachment field of \noderef{TreeSequence}; therefore its edge-node is
outside \(Q\).  Define
\[
 \sigma_n=\noderef{TwoColorComponentSwap}
       (G,\mu_n,\gamma,\alpha,\operatorname{inl}y).           \tag{5}
\]
Then
\[
                 \sigma_n(e_1)=\mu_n(e_1)=\alpha.            \tag{6}
\]

We invoke \noderef{StableChangeCasesOneTwo}, still with \(\gamma\) as the
first color.  Its component-disjointness premise is the disjointness of
\(Q\) from \(K_A\) proved above.  Its Cases 1/2 premise holds as follows:
for \(q>0\), (3) says \(\gamma\) is missing on \(T_{n,q}=A\); for \(q=0\)
non-PE, it is missing on \(T_n=A\); and for \(q=0\) PE, its construction
put it in \(\overline\mu_n(T_n)\).  Both colors are missing on the
elementary prefix \(B\) because they are missing on \(A\preccurlyeq B\).
Together with (1), the path certificate, (4), and the projections of the
minimum context, the theorem gives
\[
              \noderef{HierarchyWeaklyStable}(H,q,\sigma_n). \tag{7}
\]
\noderef{TwoColorComponentSwapProper} gives properness, and the component
swap preserves the unique uncolored edge.

We prove the pointwise missing-set equality
\[
 \overline\sigma_n(v)=\overline\mu_n(v)
                    \quad(v\in V(B)).                        \tag{8}
\]
If \(\operatorname{inl}v\notin Q\), no changed edge is incident with
\(v\).  If \(\operatorname{inl}v\in Q\), then \(v\notin V(A)\), and in
particular it is distinct from the first missing witnesses
\(v_\gamma,v_\alpha\in V(A)\).  Apply the pairwise-disjoint missing-set
clause of \noderef{ColoredVertexSetElementary} first to
\(v,v_\gamma\), using the \(\gamma\)-missing witness at \(v_\gamma\),
and then to \(v,v_\alpha\), using the \(\alpha\)-missing witness at
\(v_\alpha\).  It makes both affected colors present at \(v\).  Expanding
\noderef{MissingColors} and using
\noderef{IsProperEdgeColoring}, there is a unique incident edge of each
color.  Both edge-nodes lie in \(Q\), so (5) transposes them and keeps
both colors present; all other colors are fixed.  This proves (8), even
if \(Q\) passes through internal degree-two vertices of \(B\).
It follows in particular that the old witnesses in \(A\) keep
\(\gamma,\alpha\) missing on \(X\).                              \(\tag{9}\)

The same-good stage equalities identify the actual terminal step in the
\noderef{TashkinovSeries}.  Its PE or SE/RE
\noderef{TashkinovSeriesStep} constructor supplies the final closure.
The exact ordered-prefix and generating-coloring fields in
\noderef{ExtendedTashkinovTree} identify \(T\) with an exact ordered
prefix of that step's final series tree and identify the step coloring
with \(\mu_n\).  Unpacking the augmentation-sequence component of the
final \noderef{TashkinovClosure} gives an old route to the final series
tree, from \(T_n\) in PE and from the exact adjoined seed \(T_n+f_n\) in
SE/RE.  By \noderef{TashkinovAugmentationSequence}, this is a finite
reflexive-transitive chain of one-step augmentations under \(\mu_n\).
Use \noderef{TashkinovAugmentationSequenceRestrict} first at the exact
prefix \(T\) and then at \(X\).  Since every
\noderef{TashkinovAugmentation} increases length by one, the same finite
chain reaches a unique stage of length \(|X|\); the exact-prefix field of
\noderef{ExtendedTashkinovTree} identifies it with \(X\).  Splitting the
reflexive-transitive chain there explicitly gives both the old restricted
route to \(X\) and the old-color suffix
\[
                         X\longrightarrow T.                 \tag{11}
\]

We transport the restricted route to \(X\) by cases.  If \(q>0\), weak
stability (7) supplies its \noderef{TreeColoringInvariant} clause on
\(A\).  Every exact new edge of the old route through \(A\) is displayed
in \(A\), and every corresponding missing witness is in its predecessor
and hence in \(A\); the invariant preserves both data.  Above \(A\) the
only step before \(X\) adjoins \(e_1\), whose value is fixed by (6), and
whose old missing witness remains missing by (8).  Induction along the
finite chain therefore transports the route through \(X\), retaining the
old PE start or the exact SE/RE seed.

If \(q=0\) and the final kind is SE or RE, the nonparallel clause of
\noderef{HierarchyBase} gives \(A=T_n\).  The final
\noderef{TashkinovSeriesStep}, the exact position in (2), and the
\noderef{TreeSequenceAdjoinEdge} equality identify \(e_1=f_n\), with the
same singleton \(F_n=\{f_n\}\) and the same exact seed
\(T_n+f_n=X\).  Equation (6) retains that exact seed, so the required
route from it to \(X\) is reflexive.

Finally suppose \(q=0\) and the final kind is PE.  The normalized premise
gives an old \(\alpha\)-missing witness in \(T_n\), while the chosen
parallel-base intersection in (3) gives an old \(\gamma\)-missing witness
in \(T_n\).  The terminal ladder tree is an exact prefix of \(A\), by the
parallel clause of \noderef{HierarchyBase}, and \(A\preccurlyeq B\) by
the nested fields of \noderef{TashkinovHierarchy}.  Hence (8) preserves
both concrete witnesses under \(\sigma_n\).  Unfolding
\noderef{TreeMissingColors} turns them into
\[
 \gamma,\alpha\in\overline\sigma_n(T_n).
\]
Apply \noderef{TashkinovAugmentationSequenceComponentSwap} to the old
restricted route \(T_n\longrightarrow X\), with old coloring \(\mu_n\),
colors \(\gamma\) then \(\alpha\), component root
\(\operatorname{inl}y\), the defining swap (5), and
\(\gamma\ne\alpha\) from (3).  Its conclusion is exactly the required
\(\sigma_n\)-route from \(T_n\) to \(X\); no route-preservation claim is
inferred from \noderef{CombinedStableColoring}.

Thus in all cases we have constructed the verified route
\[
\begin{cases}
T_n\longrightarrow X,&\text{in PE},\\
T_n+f_n\longrightarrow X,&\text{in SE/RE},
\end{cases}                                                   \tag{10}
\]
under \(\sigma_n\).
No assertion about the colors of the affected suffix edges in (11) is
made.

Every \(\eta\in\noderef{HierarchyStoredColorsAt}(H,q)\) is not missing on
\(A\), whereas both \(\gamma\) and \(\alpha\) are missing there.  Hence
\(\eta\ne\gamma,\alpha\).  The component swap changes no edge to or from
color \(\eta\), and therefore
\[
 \noderef{HierarchyFirstMissingPosition}(\sigma_n,T,\eta)
 =\noderef{HierarchyFirstMissingPosition}(\mu_n,T,\eta).     \tag{12}
\]
The affected-Gamma cutoff is vacuous.  The construction (3) excludes
\(\gamma\) from every last Gamma pair.  The color \(\alpha\) is also
excluded: if \(\alpha\in
P.\mathit{counterexample}.\mathit{good}.\Gamma^q_\eta\), then storedness
of \(\eta\) makes its first missing position strictly larger than
\(|A|\).  The exact edge \(e_1=T.\mathit{edge}(|A|)\) has old color
\(\alpha\), so unfolding \noderef{HierarchySegmentColors} puts
\(\alpha\) in the last \((q,\eta)\)-segment, contradicting
\noderef{GoodHierarchy}.avoidance.                            \(\tag{13}\)

Set \(\mathcal C_0=P.\mathit{corresponding}\), and write
\(\mathcal E_0\) for the original fixed ETT and
\(\mathcal E_m=\mathcal C_0.\mathit{replacementETT}\) for the
intermediate ETT.  Apply
\noderef{StableChangeCasesOneTwoSplitTail} with replacement input
\(\mathcal C_0.\mathit{replacementInput}\), ETT \(\mathcal E_m\), base
\(P.\mathit{replacementBase}\), counterexample
\(P.\mathit{counterexample}\), and context
\(P.\mathit{minimumContext}\).  Its remaining hypotheses are the old
missing colors on \(B\), the path, disjointness, swap, and Cases 1/2 data
above, the exact split \(A\prec X\preccurlyeq B\), fixed-prefix equation
(6), verified route (10), old suffix (11), split missingness (9),
vertexwise equality (8), stored first-position equality (12), and the
vacuous cutoff from (3) and (13).  Retain the returned maximum property,
base, hierarchy, good hierarchy, and same-good data, and denote the
returned correspondence certificate by
\[
 \mathcal C_1:\quad
 \noderef{CorrespondingExtendedTashkinovTree}
    (\mathcal E_m,\sigma_n,T).
\]
Thus \(\mathcal C_1\) is relative to the intermediate ETT
\(\mathcal E_m\), not to \(\mathcal E_0\).

At \(y\), the old color \(\alpha\) is missing.  If \(\gamma\) is also
missing there, it remains so; otherwise properness gives a unique
incident \(\gamma\)-edge in \(Q\), and (5) changes it to \(\alpha\).
Thus \(\gamma\) is missing at \(y\) under \(\sigma_n\), while (8) keeps
it missing at its witness in \(A\).  The two vertices are distinct by
\(A\prec T\) and vertex injectivity in \noderef{TreeSequence}, so the
pairwise-disjoint condition in
\noderef{ColoredVertexSetElementary} fails.  These outputs and the
returned hierarchy data form the new \noderef{HierarchyCounterexample}.

We next compose correspondence back to the original ETT.  Construct
\[
 \mathcal C_{10}:\quad
 \noderef{CorrespondingExtendedTashkinovTree}
    (\mathcal E_0,\sigma_n,T)
\]
field by field.  Take its replacement input, final series, and final ETT
from \(\mathcal C_1\), transporting the series' round parameter along
\(\mathcal C_0.\mathit{rungAgreement}\).  Its rung equality is the
transitive composite
\[
 \mathcal C_1.\mathit{rungAgreement}\mathbin{;}
 \mathcal C_0.\mathit{rungAgreement}.
\]
The heterogeneous series-presentation equality of \(\mathcal C_1\)
remains the required one after this cast.

For a positive ladder index of \(\mathcal E_0\), cast first to the
corresponding index of \(\mathcal E_m\) using the inverse of
\(\mathcal C_0.\mathit{rungAgreement}\).  The new ladder tree equals the
ladder tree in \(\mathcal E_m.\mathit{series}\) by
\(\mathcal C_1.\mathit{ladderTreesAgree}\); the
\(\mathcal C_0.\mathit{seriesPresentation}\) field identifies that
dependent series with \(\mathcal C_0.\mathit{replacementSeries}\); and
\(\mathcal C_0.\mathit{ladderTreesAgree}\) identifies the latter tree
with the ladder tree of \(\mathcal E_0\).  This three-step chain is the
ladder-tree field of \(\mathcal C_{10}\).  At a full round-plus-one
index, the identical chain, using the corresponding fields of
\(\mathcal C_1\) and \(\mathcal C_0\), supplies the stored-set,
connecting-edge-set, and extension-kind equalities.  All four chains use
the dependent casts induced by the two rung equalities, so they compare
the final series directly with the original series.  Take the
final-generating-coloring, replacement-properness, and
replacement-one-uncoloredness fields from \(\mathcal C_1\), after those
same round-index casts.

It remains to construct the replacement-stability field of
\(\mathcal C_{10}\).  Fix a positive ladder index \(i\) of
\(\mathcal E_0\).  The corresponding field of \(\mathcal C_0\) is
stability from the original series coloring to the intermediate series
coloring.  Transport the field of \(\mathcal C_1\) through
\(\mathcal C_0.\mathit{seriesPresentation}\) and the ladder-tree
equality; it is stability from that same intermediate coloring to the new
series coloring on the original ladder tree.  Vertexwise equality in the
first relation, together with the raw stored-set and ladder-tree
equalities above, identifies the two computed protected
\noderef{TashkinovStoredColors} sets.

The subtype hypotheses required for transitivity in \noderef{sc2} hold
at this exact index.  The intermediate and new colorings are proper and
leave exactly \(e_0\) uncolored by the corresponding fields of
\(\mathcal C_0\) and \(\mathcal C_1\).  For the original coloring,
inspect the certified positive \noderef{TashkinovSeriesStep}; its initial
constructor cannot occur at \(i+1\).  In SE, the current-coloring equality
identifies it with the selected coloring, whose properness and
one-uncoloredness are fields of \noderef{MaximumDefectiveChoice}.  In PE
they are the current-proper and current-one-uncolored fields.  In RE the
current-coloring equality identifies it with the preceding coloring, for
which the constructor supplies previous-proper and
previous-one-uncolored.  Put the original, intermediate, and new
colorings into the subtype of \noderef{sc2}; transitivity composes the two
stable relations.  This gives the \(i\)-th replacement-stability field and
completes \(\mathcal C_{10}\).

The returned maximum property, base, hierarchy, good hierarchy, and new
counterexample remain indexed by
\(\mathcal C_1.\mathit{replacementETT}\), definitionally the final ETT
in \(\mathcal C_{10}\); both certificates also use
\(\mathcal C_1.\mathit{replacementInput}\) after the displayed cast.
Unpack the fields of \(\mathcal C_{10}\).  Its composite rung agreement
transports the lower-rung package; its composite replacement-stability at
the terminal ladder tree transports the \(2n+11\) estimate; and its
properness and one-uncoloredness fields supply those coloring hypotheses.
The returned same-good relation identifies the intermediate and final
divider counts and stages.  Together with \(P.\mathit{sameGood}\), it
transports the fewer-divider field and
\noderef{HierarchyCounterexampleMeasure}, so the original rung-local
minimality field applies.  These data give every field of the replacement
\noderef{HierarchyMinimumCounterexampleContext}, with precisely
\(\mathcal C_{10}.\mathit{replacementInput}\), its final ETT, base, and
counterexample as indices.

Finally compose the two \noderef{SameGoodHierarchy} witnesses field by
field.  In their \noderef{SameHierarchyStages} components, compose the
core equalities, divider-count equalities, heterogeneous stage and
divider-position equalities, PE equivalences, and, in the PE case,
closure equalities.  At each divider index cast along the composite count
equality, compose the stored-set equalities and then the Gamma-set
equalities for every color in that common stored set.  This gives
same-good data from the original \(C.\mathit{good}\) directly to the
final good hierarchy.  Using \(\mathcal C_{10}\), rather than
\(\mathcal C_1\), as the corresponding field, the preceding data form a
genuine replacement presentation \(P_1\) relative to the original fixed
ETT.

The Situation 7.1 condition also transports to \(P_1\).  Indeed,
\noderef{SameGoodHierarchy} gives equal divider counts and
\noderef{SameHierarchyStages}; their exact last starred stages agree.
The definition of \noderef{HierarchyPathNumber} then has the same tail
length, the same exact displayed edges and endpoints, and the same
candidate filter, so
\[
 \noderef{HierarchyPathNumber}
       (P_1.\mathit{counterexample}.\mathit{hierarchy})=1.    \tag{14}
\]
The selection of \(\gamma\), together with (8) on \(A\), also gives the
normalized positive- or zero-divider missingness for \(P_1\).

If some \(1\le i\le\ell-1\) already has a color common-missing at
\(y_i,y_\ell\) under \(P_1\), then \(i,P_1\), and that color directly
witness membership in
\noderef{HierarchyEligibleCommonMissingIndices}.  Otherwise we have the
full no-common premise required by both earlier branches.  Inspect the
finite set of exact tail positions whose \(\sigma_n\)-value is
\(\gamma\).  If it is empty, apply
\noderef{HierarchyClaimNineUnusedTailColor} to
\(P_1,\gamma\), using (14) and the transported normalization.  If it is
nonempty, choose its least position and write it as
\(|A|+j-1\).  Equation (6) and \(\gamma\ne\alpha\) show \(j\ne1\);
the tail bounds give \(2\le j\le\ell\), and minimality gives the exact
earlier-edge exclusions.  Apply
\noderef{HierarchyClaimNineFirstUseAfterSeed} to
\(P_1,\gamma,j\), again using (14).  Both branches conclude that the
eligible-index set is nonempty, as required.
\end{proof}
